\usepackage{amsmath, amssymb, amsthm, array, booktabs,xfrac, mathrsfs}
\usepackage{gensymb}
\usepackage{placeins}  % Add this to your preamble
\usepackage{float}
\usepackage[utf8]{inputenc}
\usepackage{multirow}
% to have for loop in latex
% Example:
% \foreach \x in {1,...,3}{
% \input{sol/p\x}
% }
\usepackage{mathrsfs}
\usepackage{pgffor} 
\usepackage[T1]{fontenc}
%
\usepackage[left=2.5cm, right=2.5cm, top=2cm]{geometry}
\usepackage{stackengine,scalerel}
\usepackage{commath} % using \abs
\usepackage{physics} % for partial and derivative
% Defining new command to color cell of a table
\usepackage[table]{xcolor}
    \newcommand\g{\cellcolor{green!10}}
    \newcommand\rr{\cellcolor{red!10}}
    \newcommand\yy{\cellcolor{yellow!10}}
% Algorithms
\usepackage{algorithm}
\usepackage[noend]{algpseudocode}
%
\usepackage{graphicx, amsfonts, bm, multicol, booktabs}
\usepackage[listings,many]{tcolorbox}

\usepackage{listings}
% To label row and column of a matrix
% check out :
% https://tex.stackexchange.com/questions/59517/label-rows-of-a-matrix-by-characters
\usepackage{blkarray} 
%
\usepackage{fancyhdr, lastpage,float}
\usepackage[shortlabels]{enumitem}
\usepackage{tikz,subfig}
\usepackage{epigraph,tabularx}
\usepackage{empheq}
%\usepackage{siunitx}
\usepackage{makecell} %to have a thicker hline (called xline)
\usepackage[flushleft]{threeparttable}
\usepackage{wrapfig}
\usetikzlibrary{decorations.pathreplacing,angles,quotes}
\usetikzlibrary{calc,trees,positioning,arrows,fit,shapes,calc}
% Draw cross or circle in tikz environment
\usetikzlibrary{shapes.misc}
\usetikzlibrary{patterns}
\tikzset{cross/.style={cross out, draw=black, minimum size=2*(#1-\pgflinewidth), inner sep=0pt, outer sep=0pt},
cross/.default={6pt}}    % default radius will be 6pt. 
%
% \usepackage{capt-of}
\usepackage[justification=centering, font=small,labelfont=bf]{caption}
\usepackage{eucal}
\usepackage[makeroom]{cancel}
\renewcommand{\arraystretch}{1.2}
\allowdisplaybreaks % to allow align goes to next page
% Circle
\newcommand*\circled[1]{\tikz[baseline=(char.base)]{
        \node[shape=circle,draw,inner sep=2pt] (char) {#1};}}

% x mark
\usepackage{pifont}% http://ctan.org/pkg/pifont
\newcommand{\xmark}{\text{\ding{55}}}

% Theorem 
\theoremstyle{definition}
\newtheorem{exa}{Example}
\newtheorem{nt}{Note}
\newtheorem{fact}{Fact}
% TABLE FOR INTERVAL NOTATION
\usepackage{cellspace}
\setlength\cellspacetoplimit{10pt}
\setlength\cellspacebottomlimit{4pt}
\addparagraphcolumntypes{X}
\newcolumntype{C}[1]{>{\centering\arraybackslash}m{#1}}
\newcolumntype{Y}{>{\centering\arraybackslash}X}
% Clickable Links
\usepackage{hyperref}
\hypersetup{
    colorlinks,
    citecolor=black,
    filecolor=black,
    linkcolor=black,
    urlcolor=black
}


\usepackage{amsmath}

% x and y axis style 
\usepackage{pgfplots}
\usepgfplotslibrary{polar}
\usepgflibrary{shapes.geometric}
\usetikzlibrary{calc}
\pgfplotsset{my style/.append style={axis x line=middle, axis y line=middle, xlabel={$x$}, ylabel={$y$}}}
\pgfplotsset{compat=1.15}

% NO INDENT 
\setlength\parindent{0pt}

\usepackage{minted}

\newcommand{\block}[1]{%
  \raisebox{\dimexpr(\fontcharht\font`X-1em)/2}{\rule{1em}{#1\dimexpr1em/8}}%
}
\DeclareUnicodeCharacter{2588}{\block{8}}

%%%%%%%%%%%%%%%%%%%%%%%%%%
\usepackage[thinc]{esdiff}
% examples
% \diff{f}{x} is the same as \frac{df}{dx}
% \diffp{u}{x} is the same as \frac{\partial u}{\partial x}